\documentclass[11pt,a4paper]{article}
\usepackage[T1]{fontenc}
\usepackage[utf8]{inputenc}
\usepackage[english]{babel}
\usepackage{amsmath,amssymb}
\usepackage[margin=2cm]{geometry}
\setlength{\parindent}{0pt}
\setlength{\emergencystretch}{3em}
\allowdisplaybreaks

\title{About \texttt{book\_inequalities\_with\_parameter.tex}}
\author{}
\date{}
\begin{document}
\maketitle
This note explains what the test data in \texttt{book\_inequalities\_with\_parameter.tex} is used for. That file is input for the engine in \texttt{main.c}; it is read by the test \texttt{test\_equations\_inequalities.c} in this folder.

\begin{enumerate}
\item \textbf{Purpose.} The file holds 17 true statements taken from the exercises of a school textbook on inequalities with a real parameter: each says that the inequality holds exactly in the cases the book prints as its result. The engine must prove every one of them, so the book's results are checked, not trusted.

\item \textbf{What it exercises.} The cases on the parameter (rule 18t): the comparisons that contain $x$ are read as polynomials in $x$ whose coefficients are polynomials in the parameter; the real roots of their leading coefficients, discriminants and resultants cut the line of the parameter into points and intervals, and on each of them one value of the parameter decides, with the sign table of rule 18p. Before a statement is declared proved it is also evaluated directly on a grid of pairs of rational numbers.

\item \textbf{Variety.} Second-degree inequalities whose two roots swap at a value of the parameter, a leading coefficient that vanishes ($ax^{2} + \ldots$, $(b + 1)x^{2} + \ldots$), endpoints such as $-\frac{3}{a}$ that have no value for one value of the parameter, a case where every $x$ is a solution and one where none is, and two statements with a quantifier on $x$ inside, one of them with the irrational endpoints $\pm\sqrt{7}$. Exercises 224 and 226 are not here: their printed results have $\left|2a - 1\right|$ and $\sqrt{4 - k^{2}}$ beside $x$, which the rule does not treat; their requests are in \texttt{solution\_sets\_with\_parameter.tex}.

\item \textbf{Expected result.} The line \texttt{PASS: 17 of 17 statements proved, each with a written proof}.

\item \textbf{How to run it.} From the top folder of the repository:
\begin{verbatim}
cc -std=c11 -O2 -Wall -Wextra -o /tmp/test_equations_inequalities \
    tests/equations_inequalities/test_equations_inequalities.c \
    && /tmp/test_equations_inequalities
\end{verbatim}
\end{enumerate}
\end{document}
